\section{Finite numerical closure protocol for the R1 observable kernels}
\label{sec:bhsm-kernel-closure}

After Theorems 1--5, the remaining cosmological calculation is finite.  It is
not a search for a new degree of freedom or a new phenomenological amplitude.
It is the evaluation of the already-defined reduced response kernels on the
frozen R1 branch.

\subsection{Frozen closure set}

Define
\begin{equation}
\mathfrak K_{\rm R1}
=
\left\{
\mathcal F_{D_L}(z),
\mathcal F_H(z),
\mathcal F_{D_A}(z),
\mathcal F_{f\sigma_8}(z),
\mathcal F_{\rm lens}(z)
\right\}.
\end{equation}
The branch is numerically closed only when each member of
$\mathfrak K_{\rm R1}$ is generated from one common reduced state
$X_2=(q_2,\Pi_2)^T$, one common metric response $\mathcal R_{g2}$, one common
matter response $\mathcal R_{m2}$ and one common background realization.

\subsection{No-retuning theorem}

\paragraph{Theorem 6 (frozen-kernel closure).}
Let the R1 background, physical-mode normalization, response operators and
supernova calibration be fixed.  Then no independent anisotropic nuisance
amplitude may be introduced into any member of $\mathfrak K_{\rm R1}$.
Every observable amplitude is generated by
\begin{equation}
\mathcal A_{\mathcal O}(z)
=
A_2\,\mathcal F_{\mathcal O}(z),
\end{equation}
with the same $A_2$ fixed by the single calibration theorem.

If any observable requires a second directional amplitude to match data,
that result is evidence against the pure one-mode R1 realization rather than
permission to refit the branch.

\paragraph{Proof.}
By Theorem 4, all linear observable amplitudes depend on the same unique
single-mode amplitude $A_2$.  The response operators are fixed functions of
the frozen branch.  Introducing a second directional amplitude would increase
the physical mode count or alter the frozen transfer map.  Either operation
defines a different model and therefore cannot be part of the R1 closure
calculation. \hfill$\square$

\subsection{Required numerical outputs}

A complete R1 closure record shall contain, for every requested redshift bin:

\begin{enumerate}
\item the background quantities $E(z)$, $\Omega_T(z)$ and the curvature
      scale $K/a^2$;
\item the normalized physical state $\widehat X_2(z)$;
\item the metric response amplitudes
      $\Phi_{2,\rm amp}(z)$ and $\Psi_{2,\rm amp}(z)$;
\item the matter response amplitudes relevant to density and velocity growth;
\item the five observable kernels in $\mathfrak K_{\rm R1}$;
\item the calibrated observable amplitudes obtained from the one supernova
      normalization;
\item a residual record for every defining reduced equation and transfer
      identity;
\item immutable hashes of the R1 manifest and the numerical output file.
\end{enumerate}

\subsection{Internal consistency identities}

The following identities are mandatory checks rather than optional
diagnostics:
\begin{align}
\mathcal R_{g2}
&=
-\mathcal A_g^{-1}\mathcal B_{g2},
\\
\mathcal K_2
&=
\mathcal C_2
-
\mathcal B_2^\dagger\mathcal A_g^{-1}\mathcal B_2,
\\
Y_{m,2}
&=
\mathcal R_{m2}X_2,
\\
\mathcal A_{\mathcal O}
&=
\mathcal L_{\mathcal O}\mathcal R_2 X_2.
\end{align}
Where an effective four-dimensional representation is used, its reduced
physical response must agree with these identities on the same two-dimensional
phase space.  Failure of that equivalence leaves the EFT representation open
without invalidating the BHSM-native reduction.

\subsection{Prospective-data protocol}

The R1 kernels are to be evaluated before comparison with future
high-redshift measurements.  In particular, growth and lensing kernels must
not be recalibrated to DESI, Euclid or Rubin data after those measurements are
examined.  Those datasets therefore test the forward transfer map rather than
participate in its calibration.

\subsection{Closure verdicts}

The numerical campaign has exactly three allowed verdicts:

\begin{description}
\item[R1 kernel closure.] All five kernels are produced from the common
response chain and satisfy the declared residual and provenance checks.

\item[R1 branch falsified.] The common one-mode chain is internally
inconsistent or prospective observations reject its linked amplitudes/axis.

\item[Action provenance still conditional.] The bridge kernels are numerically
well defined, but Theorem~5 has not yet been discharged on the retained BHSM
parent action.  The outputs remain BHSM-native conditional bridge
predictions.
\end{description}

No fourth verdict based on hidden retuning is permitted.
